% AY2025 力学 第2問 転記（問題のみ）
\section*{第2問〔II〕}

図2のように，半径 $a$，質量 $M$，慣性モーメント $Ma^2/2$ の均一な円盤を，水平とのなす
角 $\theta$ の粗い斜面上におき静かに手をはなすと，円盤は斜面上を滑らずに回転しながら
斜面下方に移動した。ここで，重力加速度の大きさを $g$，斜面と円盤との静止摩擦係数を
$\mu$ とする。このとき，次の問い (1)〜(3) に答えよ。

\begin{enumerate}[label=(\arabic*)]
  \item 円盤に働く摩擦力の大きさを求めよ。
  \item 円盤が滑らずに回転しながら斜面を下方に移動する条件を求めよ。
  \item 円盤の重心の加速度の大きさを求めよ。
\end{enumerate}

\begin{center}
% 図2：水平とのなす角θの斜面上に置かれた円盤（半径a）
\begin{tikzpicture}[>=stealth, scale=1.0]
  \draw (0,0) -- (5.0,0) -- (5.0,3.0) -- cycle;     % 直角三角形の斜面
  \draw (0.95,0) arc (0:31:0.95); \node at (1.3,0.25) {$\theta$};
  % 斜面上の円盤（半径0.45、斜面に接する）
  \coordinate (C) at (3.67,2.73);
  \draw (C) circle (0.45);
  \draw[->] (C) -- ++(0.33,0.31) node[above right=-2pt] {$a$};
  \node at (2.5,-0.55) {図 2};
\end{tikzpicture}
\end{center}
